\chapter*{Sobre o Open Logic Project}
\addcontentsline{toc}{chapter}{Sobre o Open Logic Project}


O \textit{Open Logic Text} é um livro-texto colaborativo de código aberto
de meta-lógica formal e métodos formais, começando em um nível intermediário
(ou seja, após um curso introdutório de lógica formal). Embora voltado a um
público não matemático (em particular, estudantes de filosofia e
ciência da computação), é rigoroso.

A cobertura de alguns tópicos atualmente incluídos pode
ainda não estar completa, e muitas seções ainda exigem revisão substancial.
Planejamos expandir o texto para cobrir mais tópicos no
futuro. Também planejamos adicionar recursos ao texto, como um
glossário, uma lista de leituras complementares, notas históricas, imagens, melhores
explicações, seções que explicam a relevância dos resultados para
filosofia, ciência da computação e matemática, e mais problemas e
exemplos. Se você encontrar um erro ou tiver uma sugestão,
\href{https://github.com/OpenLogicProject/OpenLogic/wiki/Contributing}{informe a equipe do projeto}.

O projeto opera no espírito do código aberto. O texto não só está
disponível gratuitamente; fornecemos o código-fonte LaTeX sob a
licença Creative Commons Atribuição, que dá a qualquer pessoa o direito de
baixar, usar, modificar, reorganizar, converter e redistribuir nosso
trabalho, desde que deem o devido crédito.
Consulte o site do Open Logic Project em
\href{http://openlogicproject.org/}{openlogicproject.org} para
informações adicionais.
